% Variante : corps écrit pour la v0, compilé par la v1.
%
% Seul le préambule change lors d'une migration — deux lignes. Le corps du
% document, lui, n'est pas retouché : c'est ce que garantit ocots-compat.sty,
% et c'est ce qui permet de migrer poly/ et slides/ au fil de l'eau plutôt
% qu'en bloc (~600 occurrences).
\documentclass[11pt]{ocots-book}
\usepackage[lang=fr, theme=legacy, solutions=end, math=analysis]{ocots}
\lstset{language=[LaTeX]TeX}

\title{Variante : syntaxe v0}
\author{Prénom \textsc{Nom}}
\date{\today}

\newcommand{\ocotsvariantnote}[1]{\begin{quote}\itshape Variante : #1\end{quote}}

\begin{document}

\begin{lstlisting}
\documentclass[11pt]{ocots-book}
\usepackage[lang=fr, theme=legacy, solutions=end, math=analysis]{ocots}
\end{lstlisting}
\ocotsvariantnote{deux écarts par rapport au défaut. D'abord \texttt{theme=legacy}
(le défaut est \texttt{theme=ocots}), sans rapport avec la compatibilité —
choix arbitraire hérité de la v0. Ensuite et surtout : le corps ci-dessous
emploie les \textbf{anciens noms de macros de la v0} (\texttt{mytheorem},
\texttt{myremark}, \texttt{myexample}, \texttt{myexercisecb}, \texttt{myemph},
\texttt{cblue}, \texttt{tagProblem}...) au lieu des noms actuels
(\texttt{theorem}, \texttt{remark}, \texttt{example}, \texttt{exercise},
\texttt{keyword}, \texttt{emphA}, \texttt{problemtag}...) —
\texttt{ocots-compat.sty} les fait encore fonctionner sans retouche, pour
permettre une migration au fil de l'eau.}

\tcbstartrecording

\chapter{Chapitre écrit avec les noms de la v0}
\section{Une section}

Du texte avec \myemph{une mise en valeur} et \cblue{de la couleur}.

\begin{mytheorem}{Un théorème}{v0}
    Énoncé.
\end{mytheorem}

\begin{myremark}
    Une remarque.
\end{myremark}

\begin{myexample}
    Un exemple.
\end{myexample}

\begin{proof}
    Démonstration.
\end{proof}

\begin{equation}
    \dot{x}(t) = A\, x(t) \tagProblem
\end{equation}

\begin{myexercisecb}<etiquettev0>
    Énoncé de l'exercice, avec l'ancienne syntaxe d'étiquette.
\solutioncb
    Corrigé.
\end{myexercisecb}

\begin{myexercisecb}[no solution]
    Exercice sans corrigé, ancienne clef \texttt{no solution}.
\end{myexercisecb}

\chapter{Corrections}
\tcbstoprecording
\tcbinputrecords

\end{document}
